\section{自由能把配分函数变成热力学量}
\label{sec:13-Machine-Learning/13-Statistical-Physics-and-Free-Energy/01}

你训练了一台受限玻尔兹曼机，两个 checkpoint 的重构误差几乎一样，同事问哪个模型更好，你答不上来：对数似然里挂着一个 \(\log Z\)，那是对指数多个状态求和。同一周你还在拧两个旋钮，一个是 Gibbs 采样器的温度，一个是 softmax 前面除的那个温度，文档只写了一句``温度越高越随机''。这两件事看起来毫不相干，其实是同一个量的两侧：\(\log Z\) 与温度一起出现在自由能里。本节要说清自由能是什么、为什么它比 \(Z\) 好用，以及它和变分推断里的 ELBO 是不是同一个东西。

\subsection{自由能把``能量最低''与``熵最大''写进一条式子}

设系统的微观状态是 \(x\)，能量是 \(E(x)\)，温度是 \(T>0\)。Boltzmann 分布与配分函数为

\[
p_T(x)=\frac{e^{-E(x)/T}}{Z(T)},
\qquad
Z(T)=\sum_x e^{-E(x)/T},
\]

而自由能定义为 \(F(T)=-T\log Z(T)=-T\,\logsumexp\bigl(-E/T\bigr)\)。这个定义看不出什么名堂，直到把 \(\log p_T\) 代回熵里算一遍：

\[
\begin{aligned}
S(p_T)&=-\sum_x p_T(x)\log p_T(x)
=-\sum_x p_T(x)\Bigl(-\frac{E(x)}{T}-\log Z(T)\Bigr)\\
&=\frac{\E_{p_T}[E]}{T}+\log Z(T),
\end{aligned}
\]

两边乘 \(T\) 再移项，就得到全节最该记住的一行：

\[
F(T)=\E_{p_T}[E]-T\,S(p_T).
\]

平衡态同时被两种诉求拉扯。能量项希望概率全部堆在最低的那个状态上，那样 \(\E[E]\) 最小；熵项希望概率摊得越平越好，那样 \(S\) 最大。两者方向相反，无法同时满足，自由能把它们放进一条式子里做减法，于是``该听谁的''有了定量答案。温度就是这场谈判的兑换率：降低一个单位的期望能量，值 \(1/T\) 个单位的熵。\(T\) 小的时候能量说了算，分布收缩到低能状态；\(T\) 大的时候熵说了算，分布铺向均匀。

用一个只有两个状态的系统能把这件事画出来。设状态 \(A\) 能量为 \(0\)、状态 \(B\) 能量为 \(1\)，任意分布由一个数 \(u=q(B)\in(0,1)\) 决定，于是 \(\E_q[E]=u\)、\(S(q)=-u\log u-(1-u)\log(1-u)\)。

\begin{center}
\begin{tikzpicture}[x=4cm,y=2cm,>=stealth,line join=round]
% ---- 左图：能量与熵各自的方向 ----
\draw[->] (0,0) -- (1.12,0) node[below,font=\small] {\(u\)};
\draw[->] (0,0) -- (0,1.18) node[above right,font=\small] {\(\E_q[E]\ \text{与}\ S(q)\)};
\draw[dashed,gray] (0.5,0) -- (0.5,0.72);
\draw[thick] plot[domain=0:1,samples=2] (\x,{\x});
\draw[thick,dashed] plot[domain=0.004:0.996,samples=90]
  (\x,{-\x*ln(\x)-(1-\x)*ln(1-\x)});
\node[anchor=south east,font=\small] at (0.90,0.92) {\(\E_q[E]=u\)};
\node[anchor=south,font=\small] at (0.20,0.60) {\(S(q)\)};
\node[anchor=north,font=\small] at (0.52,-0.10) {能量拉向 \(u=0\)，熵拉向 \(u=1/2\)};
% ---- 右图：两者相减，极小点随温度移动 ----
\begin{scope}[shift={(5.92cm,1.425cm)},y=1.5cm]
  \draw[->] (0,-0.95) -- (1.12,-0.95) node[below,font=\small] {\(u\)};
  \draw[->] (0,-0.95) -- (0,1.10) node[above right,font=\small] {\(F[q]=\E_q[E]-TS(q)\)};
  \draw[thick] plot[domain=0.004:0.996,samples=90]
    (\x,{\x+0.5*(\x*ln(\x)+(1-\x)*ln(1-\x))});
  \draw[thick,dashed] plot[domain=0.004:0.996,samples=90]
    (\x,{\x+1.5*(\x*ln(\x)+(1-\x)*ln(1-\x))});
  \fill (0.119,-0.063) circle (1.3pt);
  \fill (0.339,-0.622) circle (1.3pt);
  \draw[dotted] (0.119,-0.063) -- (0.119,-0.95);
  \draw[dotted] (0.339,-0.622) -- (0.339,-0.95);
  \draw[thick] (0.05,0.92) -- (0.15,0.92)
    node[anchor=west,font=\small] {低温 \(T=0.5\)};
  \draw[thick,dashed] (0.05,0.74) -- (0.15,0.74)
    node[anchor=west,font=\small] {高温 \(T=1.5\)};
  \node[anchor=north,font=\small] at (0.52,-1.02)
    {升温把极小点从 \(u^*\approx0.12\) 推到 \(0.34\)};
\end{scope}
\end{tikzpicture}
\end{center}

\emph{左：能量与熵想把分布拉向不同的地方；右：两者按温度加权相减，极小点落在拉锯的平衡处。}

极小点解出来是 \(u^*=1/(1+e^{1/T})\)，正是两状态的 Boltzmann 权重。\(T\to0\) 时 \(u^*\to0\)，系统坐死在基态；\(T\to\infty\) 时 \(u^*\to1/2\)，系统彻底摆烂成均匀分布。图右半边那两个黑点之间的位移，就是温度这个兑换率被调高之后熵多买到的份额。

\subsection{退火与 softmax 温度是同一个旋钮的两种用法}

把上面的式子换个记号，机器学习里的两处温度立刻对上号。给定 logits \(z\in\R^K\)，令 \(E(k)=-z_k\)，则 Boltzmann 分布写出来是

\[
p_T(k)=\frac{e^{z_k/T}}{\sum_{j}e^{z_j/T}}=\softmax(z/T)_k,
\]

自由能则是 \(F(T)=-T\,\logsumexp(z/T)\)。温度趋于零时 \(-F(T)\to\max_k z_k\)，softmax 退化为 \(\argmax\)；温度趋于无穷时分布趋于均匀，\(-F(T)/T\to\log K\)，剩下的全是熵。采样温度、softmax 温度、蒸馏里的温度，调的都是同一个兑换率。

退火则是让兑换率随时间变化。高温阶段熵项占优，链能翻过能量壁垒到处走；随着 \(T\) 下降，能量项逐渐夺权，分布收缩到低能区域。若降温足够慢，采样器在每个温度上都近似达到该温度的平衡分布，最终停在低能状态附近。这套做法在 MCMC 里的实现见\hyperref[sec:13-Machine-Learning/08-Sampling-and-Inference/02]{``MCMC 在无法独立采样时构造平稳链''}；它有效的前提是每一档温度都要待够久，实践中``降温足够慢''往往慢到不可接受，这是退火最常被忽略的代价。

\subsection{变分原理把不可算的 \(\log Z\) 换成一个极小值问题}

自由能真正好用的地方是它有变分刻画。对任意分布 \(q\)，仿照上面定义变分自由能

\[
F[q]=\E_q[E]-T\,S(q)=\sum_x q(x)E(x)+T\sum_x q(x)\log q(x).
\]

在被求和项里同时乘除 \(e^{-E(x)/T}\)，把 \(E(x)=-T\log\bigl(p_T(x)Z(T)\bigr)\) 代进去，两项合并后得到

\[
F[q]=T\sum_x q(x)\log\frac{q(x)}{p_T(x)}-T\log Z(T)
=T\,\KL(q\,\|\,p_T)+F(T).
\]

推导只有换元与合并同类项两步，但要对账三个条件。其一，\(Z(T)\) 必须有限，离散有限状态自动满足，连续状态需要 \(e^{-E/T}\) 可积。其二，\(q\) 的支撑要含在 \(p_T\) 的支撑里，否则 KL 取 \(+\infty\)，不等式方向不变但等号无从谈起。其三，\(\E_q[E]\) 要有限，否则左边就不是一个数。三条都成立时，由 KL 非负得 \(F[q]\ge F(T)\)，等号当且仅当 \(q=p_T\)。KL 为什么必须非负、以及它为什么不对称，见\hyperref[sec:07-Information-Theory/03-Divergences/02]{``Gibbs 不等式与 KL 的方向性''}。

这条恒等式一次给出两样东西。一是平衡分布的刻画：Boltzmann 分布不是被假设的，它是自由能泛函的唯一极小点。把能量项抽掉、只留约束，同一个极小化就退化成最大熵问题，这正是\hyperref[sec:07-Information-Theory/01-Information-and-Entropy/04]{``最大熵原理：从约束到分布''}的结构。二是可计算的上界：把 \(q\) 限制在某个便于计算的分布族里做极小化，得到的 \(F[q]\) 是 \(F(T)\) 的上界，等价于 \(\log Z\) 的下界，而差距恰好是 \(T\KL(q\|p_T)\)——近似有多差，你至少知道它差在哪个量上。配分函数本身的角色见\hyperref[sec:13-Machine-Learning/08-Sampling-and-Inference/06]{``配分函数把能量变成全局概率''}。

\subsection{物理的自由能下界就是机器学习的证据下界}

上一小节的恒等式与变分推断里的 ELBO 不是``类似''，是同一个式子换了记号。取含隐变量的模型 \(p(x,z)\)，固定观测 \(x\)，令温度 \(T=1\)，把能量定义为

\[
E(z)=-\log p(x,z).
\]

于是配分函数 \(Z=\sum_z e^{-E(z)}=\sum_z p(x,z)=p(x)\) 就是证据，自由能 \(F=-\log p(x)\) 就是负对数证据。把近似后验 \(q(z)\) 代入变分自由能：

\[
\begin{aligned}
F[q]&=\E_q\bigl[-\log p(x,z)\bigr]-S(q)
=-\Bigl(\E_q\bigl[\log p(x,z)\bigr]+S(q)\Bigr)\\
&=-\mathrm{ELBO}(q),
\end{aligned}
\]

而上一小节的恒等式 \(F[q]=\KL(q\|p_T)+F\) 逐项翻译过来就是

\[
\mathrm{ELBO}(q)=\log p(x)-\KL\bigl(q(z)\,\big\|\,p(z\given x)\bigr).
\]

对照表可以逐行核对：能量对负对数联合似然，配分函数对证据 \(p(x)\)，自由能对负对数证据，变分自由能对负 ELBO，物理里那个``近似分布与平衡分布的差距''对机器学习里那个``下界与真值的间隙''。同一个 KL，同一个方向，同一个非负性。这个对应在\hyperref[sec:13-Machine-Learning/08-Sampling-and-Inference/03]{``变分推断把后验近似变成优化''}里以概率语言展开，在\hyperref[sec:13-Machine-Learning/06-Probabilistic-Models/04]{``EM 用下界交替完成隐变量学习''}里被拆成交替优化，在\hyperref[sec:14-Deep-Learning/06-Variational-Autoencoders/01]{``VAE 把不可积后验变成可训练下界''}里被换成神经网络参数化。有了这层翻译，物理文献里的``变分自由能极小''和论文里的``最大化 ELBO''可以互相当作同一句话读。

\subsection{自由能的导数一次性生成宏观量}

自由能不只是一个标量，它是生成函数。把能量写成含外场的形式 \(E_h(x)=E_0(x)-hM(x)\)，则对 \(h\) 求导给出宏观量的期望 \(\E[M]=-\partial F/\partial h\)，再求一次导给出它的方差（乘上 \(1/T\)），也就是响应对涨落的关系；对 \(1/T\) 求导给出内能，对 \(T\) 求导给出熵。指数规模的求和被压成对一个标量函数求导，这是自由能比 \(Z\) 有用的操作性理由。

同一件事在能量模型训练里以另一副面孔出现。对数似然 \(\log p_\theta(x)=-E_\theta(x)-\log Z(\theta)\) 的梯度是

\[
\nabla_\theta\log p_\theta(x)=-\nabla_\theta E_\theta(x)+\E_{p_\theta}\bigl[\nabla_\theta E_\theta(X)\bigr],
\]

第二项那个``负相''正是 \(\log Z\) 的导数，也就是自由能的导数。训练能量模型之所以难，难点不在写不出梯度，而在负相需要模型自己的分布下的期望。对比散度、持续对比散度、平均场，都是估这一项的不同办法。开头那个``两个 checkpoint 谁更好''的问题到这里有了准确的说法：你缺的不是似然的公式，而是自由能的值或它的导数的一个可信估计。

\subsection{哪些对应是恒等式，哪些只是打个比方}

物理类比在机器学习里流行，也容易被滥用，所以要把两类说法分开摆。属于恒等式的一类，前面已经逐项核对过：配分函数就是归一化常数，就是隐变量模型的证据；自由能就是负对数证据；变分自由能就是负 ELBO；带温度的 Boltzmann 分布就是 \(\softmax(z/T)\)。这些对应可以代入验证，换记号即得，用错了会算错数。

属于比喻的一类要另行标注。把训练损失叫作``能量地形''时，除非你明确指定参数上的分布是 \(e^{-L(\theta)/T}\) 并说明 \(T\) 是什么，否则损失曲面上的``低谷''与热力学的低能态没有形式关系。把训练曲线上的突变叫``相变''时，缺少的是系综、序参量与热力学极限这三样东西。神经网络的训练是非平衡过程，自由能刻画的是平衡态，两者之间隔着动力学；密度怎样随时间流向平衡态是另一套工具的事，见\hyperref[sec:06-Random-Processes/06-Diffusion-and-SDE/02]{``Fokker--Planck 方程描述密度的流动''}。分清这两类，类比才是有用的直觉来源，而不是术语借用。

自由能的变分原理留下一个悬而未决的问题：把 \(q\) 限制到什么样的分布族里，极小化才既做得动又不至于错得太离谱。最省事的答案是假设变量之间彼此独立，代价与收益都在下一节。
